\subsection{Finite-Dimensional Riesz Representation}

Throughout this subsection, $\mathbb{F}$ is $\mathbb{R}$ or $\mathbb{C}$ and $\mathcal{V}$ is an inner-product space over $\mathbb{F}$.

\begin{theorem} \label{thm:riesz-representation-finite-dimensional} \hlt{Riesz Representation Theorem}\\
Suppose $\dim(\mathcal{V})<\infty$ and let $\varphi\in\mathcal{V}'$. Then:
\begin{enumerate}[label=(\roman*)]
\setlength{\itemsep}{0pt}
\item $(\exists\,\mathbf{u}\in\mathcal{V})(\forall\,\mathbf{v}\in\mathcal{V})\,\varphi(\mathbf{v})=\innerprod{\mathbf{v},\mathbf{u}}$;
\item $(\forall\,\mathbf{u}_1,\mathbf{u}_2\in\mathcal{V})\bigl[(\forall\,\mathbf{v}\in\mathcal{V})\,\varphi(\mathbf{v})=\innerprod{\mathbf{v},\mathbf{u}_1}=\innerprod{\mathbf{v},\mathbf{u}_2}\Longrightarrow\mathbf{u}_1=\mathbf{u}_2\bigr]$.
\end{enumerate}
\begin{proof}
\begin{enumerate}[label=(\roman*)]
\setlength{\itemsep}{0pt}
\item If $\mathcal{V}=\{\mathbf{0}\}$, then $\varphi(\mathbf{0})=0$ by linearity, so $\mathbf{u}=\mathbf{0}$ represents $\varphi$.\\
Suppose $\dim(\mathcal{V})=n\geq1$ and choose an orthonormal basis $(\mathbf{e}_1,\ldots,\mathbf{e}_n)$ by Corollary \ref{cor:finite-dimensional-inner-product-space-has-orthonormal-basis}. Set $\mathbf{u}\defeq\sum\limits_{j=1}^{n}\overline{\varphi(\mathbf{e}_j)}\mathbf{e}_j$.\\
Fix $\mathbf{v}\in\mathcal{V}$. By Theorem \ref{thm:orthonormal-basis-coordinates}(i), linearity of $\varphi$, and Proposition \ref{prop:basic-inner-product-identities}(iii)--(iv), $\varphi(\mathbf{v})=\sum\limits_{j=1}^{n}\innerprod{\mathbf{v},\mathbf{e}_j}\varphi(\mathbf{e}_j)=\innerprod{\mathbf{v},\mathbf{u}}$.
\item Suppose $(\forall\,\mathbf{v}\in\mathcal{V})\,\varphi(\mathbf{v})=\innerprod{\mathbf{v},\mathbf{u}_1}=\innerprod{\mathbf{v},\mathbf{u}_2}$. Then $\mathbf{u}_1=\mathbf{u}_2$ by Lemma \ref{lem:inner-products-separate-vectors}(i).
\end{enumerate}
\hlt{Structural Commentary}: An orthonormal expansion constructs the representative; inner-product separation proves uniqueness.\\
\hlt{Key Insight}: Functional values become conjugated coordinates of one vector.
\end{proof}
\end{theorem}

\begin{definition} \label{def:riesz-vector} \hlt{Riesz Vector}\\
If $\dim(\mathcal{V})<\infty$ and $\varphi\in\mathcal{V}'$, the unique $\mathbf{u}\in\mathcal{V}$ satisfying $(\forall\,\mathbf{v}\in\mathcal{V})\,\varphi(\mathbf{v})=\innerprod{\mathbf{v},\mathbf{u}}$ is the \hlt{Riesz vector} representing $\varphi$.
\end{definition}

\begin{theorem} \label{thm:riesz-map-conjugate-isomorphism} \hlt{Riesz Map}\\
Suppose $\dim(\mathcal{V})<\infty$ and define $\tau:\mathcal{V}\to\mathcal{V}'$ by $\tau(\mathbf{u})\defeq\innerprod{\cdot,\mathbf{u}}$. Then:
\begin{enumerate}[label=(\roman*)]
\setlength{\itemsep}{0pt}
\item $\tau$ is conjugate linear: $(\forall\,a\in\mathbb{F})(\forall\,\mathbf{u},\mathbf{w}\in\mathcal{V})\,\tau(a\mathbf{u}+\mathbf{w})=\overline{a}\,\tau(\mathbf{u})+\tau(\mathbf{w})$;
\item $\tau$ is bijective;
\item $R\defeq\tau^{-1}$ is conjugate linear: $(\forall\,a\in\mathbb{F})(\forall\,\varphi,\psi\in\mathcal{V}')\,R(a\varphi+\psi)=\overline{a}\,R(\varphi)+R(\psi)$;
\item if $\mathbb{F}=\mathbb{R}$, then $\tau$ and $R$ are linear.
\end{enumerate}
\begin{proof}
\begin{enumerate}[label=(\roman*)]
\setlength{\itemsep}{0pt}
\item $(\forall\,\mathbf{u}\in\mathcal{V})\,\tau(\mathbf{u})\in\mathcal{V}'$ by Proposition \ref{prop:basic-inner-product-identities}(i). Fix $a\in\mathbb{F}$, $\mathbf{u},\mathbf{w},\mathbf{v}\in\mathcal{V}$. Then $\tau(a\mathbf{u}+\mathbf{w})(\mathbf{v})=\overline{a}\,\tau(\mathbf{u})(\mathbf{v})+\tau(\mathbf{w})(\mathbf{v})$ by Proposition \ref{prop:basic-inner-product-identities}(iii)--(iv), so $\tau(a\mathbf{u}+\mathbf{w})=\overline{a}\,\tau(\mathbf{u})+\tau(\mathbf{w})$ by equality of functions.
\item If $\tau(\mathbf{u})=\tau(\mathbf{w})$, then $\norm{\mathbf{u}-\mathbf{w}}^2=\tau(\mathbf{u}-\mathbf{w})(\mathbf{u}-\mathbf{w})=0$ by (i), so $\mathbf{u}=\mathbf{w}$ by definiteness.\\
Thus $\tau$ is injective. For $\varphi\in\mathcal{V}'$, $(\exists\,\mathbf{u}\in\mathcal{V})\,\tau(\mathbf{u})=\varphi$ by Theorem \ref{thm:riesz-representation-finite-dimensional}(i), so $\tau$ is surjective.
\item Fix $a\in\mathbb{F}$ and $\varphi,\psi\in\mathcal{V}'$. By (i) and the definition of $R$, $\tau\bigl(\overline{a}\,R(\varphi)+R(\psi)\bigr)=a\varphi+\psi=\tau\bigl(R(a\varphi+\psi)\bigr)$.\\
Hence $R(a\varphi+\psi)=\overline{a}\,R(\varphi)+R(\psi)$ by injectivity from (ii).
\item If $\mathbb{F}=\mathbb{R}$, then $\overline{a}=a$ in (i) and (iii), so both maps are linear.
\end{enumerate}
\hlt{Structural Commentary}: Conjugate-linearity, separation, and representation establish the bijection and its inverse law.\\
\hlt{Key Insight}: Conjugating the representing slot conjugates scalar multiplication.
\end{proof}
\end{theorem}

Figure \ref{fig:riesz-representation-vector-functional} depicts Theorem \ref{thm:riesz-map-conjugate-isomorphism}: the vector $\mathbf{u}$ corresponds to the functional $\mathbf{v}\mapsto\innerprod{\mathbf{v},\mathbf{u}}$.

\begin{figure}[H]
  \centering
  \begin{tikzpicture}[
  >=Stealth,
  every node/.style={font=\small},
  space/.style={draw=gray!45, rounded corners=3pt, minimum width=3.7cm, minimum height=2.5cm},
  dual/.style={draw=gray!45, rounded corners=3pt, minimum width=3.7cm, minimum height=2.5cm, fill=gray!5},
  vec/.style={->, blue!75!black, line width=0.9pt},
  map/.style={->, gray!70, line width=0.75pt},
  test/.style={->, orange!80!black, line width=0.7pt}
]
  \node[space] (V) at (0,0) {};
  \node[anchor=north west] at (-1.65,1.1) {$\mathcal{V}$};
  \coordinate (O) at (-0.95,-0.65);
  \coordinate (u) at (0.75,0.65);
  \coordinate (v) at (0.8,-0.15);
  \draw[vec] (O) -- (u) node[midway, above left] {$\mathbf{u}$};
  \draw[test] (O) -- (v) node[midway, below] {$\mathbf{v}$};

  \node[dual] (Vdual) at (5.3,0) {};
  \node[anchor=north west] at (3.65,1.1) {$\mathcal{V}'$};
  \node[align=center] at (5.3,0.25) {$\varphi_{\mathbf{u}}$};
  \node[align=center] at (5.3,-0.45) {$\varphi_{\mathbf{u}}(\mathbf{v})=\innerprod{\mathbf{v},\mathbf{u}}$};

  \draw[map] (1.95,0.25) -- (3.35,0.25) node[midway, above] {Riesz};
  \draw[test] (1.0,-0.15) .. controls (2.5,-1.55) and (4.0,-1.55) .. (5.0,-0.85)
    node[midway, below] {evaluation};
\end{tikzpicture}

  \caption{The map $\tau$ sends each vector to the functional it represents.}
  \label{fig:riesz-representation-vector-functional}
\end{figure}

\begin{perspective} \label{per:riesz-representation-coordinates} \hlt{Riesz as Coordinate Duality}\\
For an orthonormal basis $(\mathbf{e}_1,\ldots,\mathbf{e}_n)$, Theorem \ref{thm:riesz-representation-finite-dimensional}(i) constructs the Riesz vector of $\varphi\in\mathcal{V}'$ as $\sum\limits_{j=1}^{n}\overline{\varphi(\mathbf{e}_j)}\mathbf{e}_j$.\\
Although this formula uses the basis, uniqueness in part (ii) makes its value basis-independent.
\end{perspective}
